\chapter{认知对象的结构化表示：语言、视觉与关系}
\label{chp:semantic_anatomy}

环境记录进入系统之后，还需要被组织成可以用于判断、推理和交流的认知对象。我们用这个名称指系统内具有身份、属性和关系的表示单位，不预设它与主观体验等同，也不要求它复制外部对象的全部细节。一个句子、一幅场景图和一条论证链都可以包含这些单位，但它们的关系规则不同。语言中的施事角色不能直接当作图像中的左侧位置，逻辑蕴含也不能直接当作事件因果。本章分别说明语言、视觉和推理中的结构内容绑定，再讨论它们怎样通过明确接口协作。我们把理解表述为一种任务能力：系统能够在语境变化下维持对象与关系，回答有证据支持的问题，并依据新信息修订原有解释。这样，认知对象的有效性可以接受检查，而不依赖向量外观、几何修辞或输出流畅程度来证明。

\section{语言表示：表达顺序、事件角色与语境消歧}
\label{sec:language_structure}

语言输入以字符、词项或其他符号序列出现，但序列顺序并不足以表达所有关系。“猫追鱼”和“鱼追猫”具有相同的对象词与动作词，事件角色却相反。如果系统只保存词项出现次数，就无法回答谁是追逐者。我们因此把表面序列、句法依赖和事件角色作为不同层次的记录。它们可以相互支持，也可能彼此不一致。错误输入、歧义表达和省略现象都要求系统保留候选解释，而不是从词序直接跳到唯一事件。

设事件 $e$ 具有动作类型和若干角色槽位，我们用 $R(e,\mathrm{actor},i)$ 与 $R(e,\mathrm{target},j)$ 记录施事和目标。这些关系包含角色类型，交换对象身份会改变事件内容。若用上一篇的张量绑定形式，可写为 $T_e=\sum_r v_r q_{e,r}^{\mathsf T}$，其中 $v_r$ 是角色编码，$q_{e,r}$ 是角色内容。角色编码必须满足读出所要求的可区分条件；否则，施事与目标可能在压缩后混淆。公式提供一种实现，但事件角色的语义来自建模规范和验证任务，不能从张量乘法本身产生。

词义需要语境。一个词项可能表示物体、动作或某个专业概念；相同词项在不同语境下的读出不应被强行固定。我们将词项的基础编码与上下文更新后的对象表示区分，要求上下文更新保存与任务相关的身份信息。例如，“把它放在桌上”中的代词需要指向已有对象；若前文出现多个候选，“它”的归属应由对话条件和任务证据决定。系统可以请求或寻找更多证据，也可以明确保留分支，而不应把最高相似度候选当成已证实身份。

语法可接受与事实正确是两个不同检查。一个句子可以结构完整却与环境记录冲突，也可以表达想象情境而没有声称其为现实。我们在认知对象中保留表达模式：事实陈述、条件假设、引用内容与反事实推演不能混为一谈。对“假如苹果是蓝色的”进行推理，不要求系统把现实苹果的颜色改成蓝色；它要求建立受条件约束的分支表示。分支结束后，事实状态仍应由原始证据支持。这种作用域管理是理解复杂语言的必要组成部分。

语义冲突也不是简单的向量相反。两个描述是否冲突，取决于对象身份、时间、属性定义和语境。苹果在昨天为青色、今天变为黄色，未必矛盾；两个来源在同一条件下给出互斥类别，才构成需要解释的冲突。我们把冲突检查落实到这些字段，并记录当前的解决依据。仅在编码空间中计算夹角，可能有助于发现候选冲突，却不能代替类型、作用域和证据条件的检查。

语言理解的评估应覆盖关系干预。我们可以保持词项不变而交换角色，检查回答是否随角色改变；可以改写表达而保留事件，检查关键读出是否保持；还可以加入无关句子，检查原对象身份是否稳定。三个试验分别检查关系敏感性、表达稳健性和干扰控制。单一问答正确率很难区分这几种能力。尤其当系统记住常见搭配时，它可能回答原句，却无法处理角色互换后的句子。

在持续对话中，新输入不仅添加内容，还可能修订旧对象。用户说明“我刚才说的苹果其实是模型道具”，系统应调整对象类别、可用属性和相关推理，而不只是追加一句备注。我们保存修订事件及其影响范围，使仍依赖旧解释的计划进入重新检查。语言表示由此成为能够维护的任务状态，而非输入文本的静态摘要。HSF-HD 在这一层提供一般的绑定、作用域与修订要求，具体语言模型只是完成这些操作的一种实现选择。

\section{视觉表示：区域、对象身份与场景关系}
\label{sec:vision_structure}

视觉输入首先是由采样装置产生的数组。对象、轮廓和空间关系需要从数组中估计，而不是直接写在像素里。同一个对象可以因视角改变而呈现不同区域，两个不同对象也可能具有相似外观。我们将区域记录、对象记录和场景关系分开：区域指向当前图像位置，对象维持跨时刻身份，关系则说明对象之间的相对组织。这样，移动对象不必改变其身份，摄像机移动也不必被解释成所有对象都改变了位置。

区域记录可以包括边界框、掩码和可见性标记。对象属性则可以包含类别候选、外观编码与估计尺度。边界框通常只是一种位置摘要，不能精确代表对象占据的全部区域；掩码表达更细的像素归属，但仍可能受到遮挡影响。我们不把检测输出当成无误的对象定义，而是保存算法版本与支持证据。当任务要求接触位置时，分类准确并不足够，还需要相应的几何估计与误差范围。

场景关系需要声明坐标系与尺度。“苹果在盘的左边”可能依据图像坐标，也可能依据操作者视角；更换视角后，同一句关系可能改变。包含、接触与遮挡又具有不同含义。遮挡关系主要表达某个观测视角下的可见性顺序，不能自动推出物体之间的接触关系。我们为关系添加参照系与观测条件，使同一场景中不同来源的关系能够被比较，而不是仅凭相同标签合并。

对象追踪使用跨帧证据维持身份。若两个外观相似的苹果交叉移动，系统可能暂时无法确定对应关系。我们可以把匹配表示为带约束的候选分配，要求一个观测候选不同时归属两个互斥对象；遇到遮挡、合并或分裂时，则需要显式事件更新。只有在证据足够时，才把候选匹配提交为新的身份关系。身份错误会传播到历史轨迹与动作计划，因此它应拥有独立于类别置信度的诊断指标。

结构与外观可以相互帮助，但这种帮助不能替代证据。例如，盘的边缘可能帮助解释被遮挡苹果的轮廓；同时，苹果的形状也能帮助判断哪条边属于盘。系统在联合估计中可以使用先验形状，但必须区分直接可见区域与补全区域。补全结果可能对预测有用，却不是新增观测。若把推断出的边界标记为已观测，后续系统就可能高估确定性，并在新的视角出现时拒绝合理修订。

我们可以通过控制结构与属性的变化检查表示分工。保持位置关系而改变颜色，适合检查对象定位是否依赖不必要的外观捷径；保持外观而交换对象位置，适合检查场景关系是否真正编码。评价这些试验时，应同时考察对象身份、关系读出和任务结果。仅凭重建图像相似，无法证明内部状态保留了所有关键关系，因为重建目标可能更关注像素而非角色和对象连续性。

视觉与语言的对齐需要对应对象，而不只是对应整体句子与整体图像。“盘里有两个苹果”要求语言中的数量、包含关系和对象类别都能在视觉记录中找到支持。若某个苹果被遮挡，系统应说明数量是直接数得、由历史追踪维持还是由模型推断。跨模态接口由此不仅传递编码，还传递证据条件。对于 AgentOS，这使视觉工具的结果能够参与可回溯的任务状态，而不会仅以一段自然语言描述取代所有原始依据。

\begin{figure}[htbp]
\centering
\includegraphics[width=0.82\linewidth,height=0.54\textheight,keepaspectratio]{image/nureal_apple.png}
\caption{苹果的内部表示示意：对象、属性与关系属于系统组织的任务状态，其有效性需由观测和读出分别检验。}
\end{figure}

\section{推理表示：证明依赖、可信程度与因果条件}
\label{sec:logic_structure}

推理需要在对象之间建立可检查的依赖。我们首先区分形式蕴含、统计相关和环境因果。形式蕴含要求结论依据规定推导规则从前提产生；统计相关描述变量在某个分布下共同变化；因果关系还涉及干预以及对干预后结果的预测。这三种关系可能在同一任务中相互协作，却不能用同一条无类型边代替。一个图上存在连通路径，也不足以说明沿路推导有效，因为路径中的关系类型和条件可能不兼容。

以命题逻辑中的一次推导为例，前提集合包含 $A$ 和 $A\Rightarrow B$，使用规定规则得到 $B$。我们记录结论节点、前提节点及所用规则，形成一个证明事件。检查器需要验证该事件的输入类型、作用域和规则实例，而不是只检查节点相邻。若前提位于一个暂时假设分支，结论也继承该分支条件；离开分支时，不能把结论直接当成无条件事实。证明依赖图因此既记录推导顺序，也记录成立条件。

命题真值与系统可信程度不能混淆。某个命题在指定语义下为真，并不意味着系统已经有充足证据知道它为真；反过来，高置信度判断也可能错误。我们分别保存形式检查结果与经验支持。概率推断要求联合分布或相应的条件独立假设，不能把两条前提置信度随意相乘来获得结论概率。若缺少这些条件，系统应使用有明确意义的评分或区间，而不把结果包装成精确概率。

依赖记录使撤回成为可计算操作。若前提 $A$ 的证据失效，所有仅依赖该证据的结论需要重新评估；若某个结论还有独立证明，则不能简单一并删除。我们为结论保存多个支持集合，在修订时检查是否仍存在有效支持。这个机制比一句“更新信念”更具体，因为它规定了什么发生变化、哪些结果受到影响，以及为何某些结果仍被保留。复杂任务中的一致性维护依赖这样的结构。

因果推理还需要区别观察条件与干预条件。看到某变量与另一变量共同变化，只能提供相关证据；主动改变前者后，后者是否随之变化，取决于系统机制和其他条件。我们可以用类型化接口分别记录观测分布与干预预测，并要求模型声明它们之间的连接假设。对无干预能力的任务，也可以提出因果候选，但应保留其证据限制。语言表达中的“因为”不是自动成立的因果证明。

寻找证明与检查给定证明也是不同计算任务。某些问题中，给定见证可以较快检查，而寻找见证需要大量搜索；这种差异依赖问题类型和表示方式，不能被宣称为所有推理的普遍成本定律。我们在任务记录中分别统计候选生成成本、检查成本和失败回溯成本。生成模型可以提出辅助条件或新候选，检查器则判断候选是否满足规则。创造性探索与严格验证因而可以协作，无须把探索禁止后才能获得可靠推理。

推理的状态更新包括增加前提、提交已验证步骤、保留未验证候选与撤回失效结论。每种更新都有不同条件。我们不把所有认知内容压成一个持续增大的真值列表，而是维护具有状态和来源的论证对象。对于 AgentOS，一个规划建议、一次工具结果和一条经检查的推导可以共享记录接口，却需要不同的有效性标记。HSF-HD 所提供的一般性在于规定依赖与更新的组织方式，而不是把所有知识都化约成同一种推理规则。

\section{跨模态协作：共同身份、接口契约与可修订理解}

语言、视觉和推理可以协作，是因为它们能够建立任务相关的对应关系，而不是因为它们天然拥有相同编码。我们首先定义共同身份接口：语言中的“这个苹果”、视觉追踪中的目标对象和规划中的操作对象，需要指向同一个任务实体或明确的候选集合。身份接口不要求三个模块保存同一个向量，但要求它们能够解释何种证据支持对应，以及对应失败时怎样修订。

接口契约规定输入类型、输出类型、有效条件与错误状态。视觉模块可以输出区域候选和对象属性，语言模块可以输出事件角色与指代候选，推理模块可以输出检查结果与依赖关系。接收方不能凭字段名字相似便认定语义一致。例如，视觉置信度可能是分类评分，语言置信度可能是生成评分，两者没有共同校准依据时不应直接平均。我们可以通过统一任务上的校准或选择规则组合它们，但必须说明组合目标。

在数值消息传递中，设对象 $i$ 的状态为 $x_i^k\in\mathbb R^d$，关系类型为 $r$，更新可以采用
\begin{equation}
x_i^{k+1}=f\left(x_i^k,\sum_r\sum_{j\in\mathcal N_r(i)}A_r x_j^k,u_i^k\right).
\end{equation}
矩阵 $A_r$ 把对应消息映射到共同消息空间，加和因此具有明确类型；函数 $f$ 再产生下一状态。这个机制能够表达关系敏感更新，却不保证自动完成逻辑证明或对象追踪。关系是否被正确建立、消息是否保留关键角色以及读出是否满足任务条件，仍需要独立验证。数值运算的可微性不等于语义正确性。

理解可以通过一组任务能力得到操作化描述。系统需要回答直接问题、处理等义改写、识别角色交换、解释来源冲突，并在新证据出现时修订判断。这些能力共同构成一个评估范围，任何单项成功都不必然推出其余能力。我们还检查更新后的状态是否保持旧任务中仍然有效的关系，避免新信息覆盖对象身份或作用域。持续任务中的理解，因此不仅表现为当前输出，还表现为修订过程是否保存了必要连续性。

以“拿起盘中较大的苹果”为例，语言解析需要给出操作类型、目标限定与比较属性；视觉处理需要定位盘、建立包含候选并估计大小；规划需要选择可行抓取动作。若两个苹果大小相近且误差区间重叠，系统应保留选择不确定性，必要时获取新的视角。若用户随后说明“我指的是左边那个”，这条输入修订目标限定，但不必改变两个苹果原有的视觉身份。各模块通过共同实体和修订事件协调，才不会把新增限定误当成新增对象。

跨模态冲突不应总由某个模块固定裁决。语言可能提供视觉无法看到的属性，视觉也可能纠正语言中的旧描述，形式检查则可能揭示两者联合时的作用域错误。我们根据证据类型、时间和任务要求选择裁决方式，并保存无法解决的冲突。系统可以暂停依赖该冲突的行动，同时继续与其无关的工作。这种局部阻断比让一个不确定判断污染全部任务状态更有利于维持持续运行。

我们最终把认知对象看作可组织、可读出和可修订的计算状态。它们的内容可以来自多种输入，其关系必须具有明确含义，其更新需要保留证据与依赖。这个定义为一般智能理论提供了共同接口，同时保留不同模态和推理机制的差异。AgentOS 可以在这一接口上实现工具协作、任务记忆与执行记录；它的具体模块并不是 HSF-HD 一般理论的前提，而是对一般状态组织要求的一种工程实现。
